% =========================================================================
% English translation (generated by AI using Claude Opus 4.8 (Max effort)).
% This file was translated from Hungarian to English by AI.
% DISCLAIMER: Please review against the original Hungarian .tex files, before relying on it!
% SOURCE: 2. Differenciál- és integrálszámítás.tex
% =========================================================================
\documentclass[margin=0px]{article}

\usepackage{listings}
\usepackage[utf8]{inputenc}
\usepackage{graphicx}
\usepackage{float}
\usepackage[a4paper, margin=0.7in]{geometry}
\usepackage{subcaption}
\usepackage{amsthm}
\usepackage{amssymb}
\usepackage{amsmath}
\usepackage{fancyhdr}
\usepackage{setspace}

\onehalfspacing

\renewcommand{\figurename}{Figure}
\newenvironment{tetel}[1]{\paragraph{#1 \\}}{}
\newcommand{\R}{\mathbb{R}}

\pagestyle{fancy}
\lhead{\it{CS BSc Final Exam Topics}}
\rhead{2. Differential and integral calculus}

\title{\textbf{{\Large ELTE IK - Computer Science BSc} \vspace{0.2cm} \\ {\huge Final Exam Topics}} \vspace{0.3cm} \\ 2. Differential and integral calculus}
\author{}
\date{}

\begin{document}
\maketitle

\begin{center}
\fbox{\parbox{0.92\textwidth}{\small \textit{Note: This document was translated from Hungarian to English by AI. Please verify the information against the original Hungarian source before relying on it.}}}
\end{center}

\begin{tetel}{Differential and integral calculus}
    Differentiability of functions. Partial derivative, total derivative. Gradient, Jacobian matrix. Function analysis, extrema. Riemann integral, integration by parts, integration by substitution. Newton--Leibniz formula.
\end{tetel}

\section{Differentiability of functions}
\begin{description}
    \item[Differentiability] \hfill \\
        $ 1 \leq n, m \in \mathbb{N}, \quad 1 \leq p,q \leq +\infty,$

        $ (\R^n, \lVert . \rVert_p)$ and $ (\R^m, \lVert . \rVert_q)$ normed spaces

        $ f \in \R^n \rightarrow \R^m, \ a \in intD_f $

        The function $f$ is differentiable at the point $a$ ($ f \in D\{a\}$), if\\
        there exists a bounded linear map $ L \in L(\R^n, \R^m)$ and a function $ \eta \in \R^n \rightarrow \R^m $ such that:
        \begin{align*}
            f(a+h)-f(a) = L(h)+\eta(h)\cdot \lVert h\rVert_p \quad ( h \in \R^n, a+h \in D_f)
        \end{align*}
        where
        \begin{align*}
            \eta(h) \longrightarrow 0 \quad (\lVert h\rVert_p \rightarrow 0)
        \end{align*}

        In other words:

        \[ \dfrac{f(a+h)-f(a)-L(h)}{\lVert h\rVert_p} \longrightarrow 0 \quad (\lVert h\rVert_p \rightarrow 0)\]
        \\\\
        If $\forall a\in intD_f : f \in D\{a\} $, then $f$ is differentiable ($f \in D$)
        \\\\
        \textit{
            Remark:\\
            A continuous map between the normed spaces $ (X, \lVert.\rVert_\bigstar )$, $ (X, \lVert.\rVert_\heartsuit )$ over the field $ \mathbb{K}$ is a bounded linear map if
            \begin{itemize}
                \item it is linear, that is,
                      \begin{align*}
                          f(x+\lambda y) = f(x) + \lambda f(y) \quad (x,y \in X, \lambda \in \mathbb{K})
                      \end{align*}
                \item it is bounded, that is,
                      \begin{align*}
                          \exists M \geq 0 : \lVert f(x) \rVert_\heartsuit \leq M\lVert x \rVert_\bigstar \quad (x \in X)
                      \end{align*}
            \end{itemize}
        }
    \item[Derivative] \hfill \\
        the function $f\in\R^n\rightarrow\R^m $ is differentiable at a point $a\in intD_f$ \\
        $ \Rightarrow \exists! L\in L(\R^n,\R^m)$ bounded linear map

        This uniquely existing bounded linear map $L \in L(\R^n,\R^m) $ is called the derivative of the function $f$ at the point $a$, and is denoted by the symbol $f'(a)$.
\end{description}

\section{Partial derivative, total derivative}
\subsection{Partial derivative}
\begin{description}
    \item[Definition] \hfill \\
        Consider the function $ h \in \R^n \rightarrow \R $ and the vector $a = (a_1, ... , a_n) \in D_h $. Let
        \begin{align*}
            D_{h,i}^{(a)} := \{t \in \R : (a_1, ..., a_{i-1}, t, a_{i+1}, ..., a_n) \in D_h\} \quad (i = 1,...,n)
        \end{align*}
        And let:
        \begin{align*}
            h_{a,i} : D_{h,i}^{(a)} \rightarrow \R, \quad \textrm{ for which: } \quad h_{a,i}(t) := h(a_1, ..., a_{i-1}, t, a_{i+1}, ..., a_n) \quad (t\in	D_{h,i}^{(a)} )
        \end{align*}
        The partial functions $h_{a,i}$ are all single-variable real functions ($h_{a,i} \in \R \rightarrow\R $)

        The function $ h $ is partially differentiable at $a$ with respect to the $i$-th variable if $ h_{a,i} \in D\{a_i\} $. Then:
        \begin{align*}
            \partial_ih(a) := h'_{a,i}(a_i)
        \end{align*}
        the real number is called the partial derivative of the function $h$ at $a$ with respect to the $i$-th variable.
    \item[Partial derivative function] \hfill \\
        Suppose that for the previous function $h$:
        \begin{align*}
            D_{h,i} := \{a \in D_h : \textrm{the partial derivative } \partial_ih(a) \textrm{ exists}\} \neq \emptyset \qquad (i=1,...,n)
        \end{align*}
        Then the function
        \begin{align*}
            x \mapsto \partial_ih(x) \quad (x\in D_{h,i})
        \end{align*}
        is called the partial derivative function of the function $h$ with respect to the $i$-th variable, and is denoted by the symbol $\partial_ih$.
    \item[Differentiability and partial differentiability] \hfill
        \begin{itemize}
            \item Differentiability $\Rightarrow$ partial differentiability \\
                  $ 1 \leq n \in \mathbb{N}, \ h \in \R^n \rightarrow \R, $ and $h\in D\{a\} \ (a \in D_h) $ \\
                  $\Rightarrow \forall i = 1,...,n$ : the function $h$ is partially differentiable at the point $a$ with respect to the $i$-th variable, and
                  \begin{align*}
                      \textrm{grad}h(a) = (\partial_1h(a), ..., \partial_nh(a))
                  \end{align*}

            \item Differentiability $\Leftarrow$ partial differentiability \\
                  $ 1 \leq n \in \mathbb{N}, \ h \in \R^n \rightarrow \R, $ and $ a \in intD_h $ \\
                  For some $i = 1,...,n$:
                  \begin{itemize}
                      \item At any point $ x \in K_r(a) $ ($r > 0$ suitable) the partial derivatives $ \partial_jh(x)$ exist $ (i \neq j = 1,...,n)$ and these are continuous
                      \item $\exists \ \partial_ih(a)$ partial derivative
                  \end{itemize}
                  $\Rightarrow h\in D(a)$
        \end{itemize}
\end{description}
\subsection{Total derivative}
\rm The function $ f: \R^p \to \R^q $ is said to be {\it (totally) differentiable\,} at the point $a\in D(f)$,
if there exists a linear map $ L: \R^p \to \R^q $ such that
$$
    \lim_{x\to a}\frac{f(x)-f(a)-L(x-a)}{\|x-a\|}=0.
$$
The linear map $L$ is called the {\it differential} of the function $f$ at the point $a$, and the matrix $M_L$ of the map $L$ is called
the {\it differential quotient or Jacobian matrix\,} of $f$ at $a$. Notation: $L=Df(a)$, and $M_L=f'(a)$.

For real $f$ ($q=1$) the Jacobian matrix $f'(a)$ is of type $1\times p$, that is, a $p$-dimensional row vector. In this case, instead of the
Jacobian matrix at $a$, the name {\it gradient {\rm or} gradient vector\,} at $a$ and the notation $f'(a)=\operatorname{grad} f(a)$
are also used.

\section{Jacobian matrix, gradient}
\begin{description}
    \item[Jacobian matrix] \hfill \\
        For the bounded linear map $ L := f'(a) $ appearing above, $ \exists!A\in\R^{m\times n}$ matrix for which:
        \begin{align*}
            L(x) = Ax \quad (x\in\R^n)
        \end{align*}
        Therefore:
        $ f'(a) := A $\\
        is the derivative or derivative matrix of the function $f$ at $a$, also known as its Jacobian matrix.
    \item[Gradient] \hfill \\
        For $m = 1$: $ f \in \R^n \rightarrow \R $
        \begin{align*}
            \textrm{grad}f(a) := f'(a) \in \R^{1\times n} \approx \R^n
        \end{align*}
        So in this case the Jacobian matrix $f'(a)$ can be regarded as a vector in $\R^n$, which is called the gradient of the function $f$ at $a$.
        \\\\
        If $ D := \{a \in D_f : f \in D\{a\} \} $, then the function
        \begin{align*}
            x \mapsto \textrm{grad}f(x) \in \R^n \quad (x \in D)
        \end{align*}
        is called the gradient of the function $f$, and is denoted grad$f \in \R^n \rightarrow \R^n$.
    \item[Gradient as a row of the Jacobian matrix] \hfill \\
        Let $1 \leq n, m \in \mathbb{N}$. The function $f = (f_1, ..., f_m) \in \R^n \rightarrow \R^m$ is differentiable at the point $ a \in intD_f$ if and only if for every $ i = 1, ..., m $ the
        coordinate function $f_i \in \R^n \rightarrow \R$ is differentiable at $a$.

        If $f \in D\{a\}$, then the Jacobian matrix $f'(a)$ has the following form:
        \begin{align*}
            f'(a) =
            \begin{bmatrix}
                \textrm{grad}f_{1}(a) \\
                \textrm{grad}f_{2}(a) \\
                \vdots                \\
                \textrm{grad}f_{m}(a) \\
            \end{bmatrix}
        \end{align*}
\end{description}

\section{Extrema, function analysis}
\subsection{Extremum}
$ f \in \R \rightarrow \R, a \in D_f$
\begin{description}
    \item[Local maximum] \hfill \\
        $f$ has a local maximum at $a$ if for a suitable $ r > 0$: \\
        $ f(x) \leq f(a) \qquad (x\in D_f, |x-a|<r)$
    \item[Local minimum] \hfill \\
        $f$ has a local minimum at $a$ if for a suitable $ r > 0$: \\
        $ f(x) \geq f(a) \qquad (x\in D_f, |x-a|<r)$
    \item[Absolute maximum] \hfill \\
        $f$ has an absolute maximum at $a$ if: \\
        $ f(x) \leq f(a) \qquad (x\in D_f)$
    \item[Absolute minimum] \hfill \\
        $f$ has an absolute minimum at $a$ if: \\
        $ f(x) \geq f(a) \qquad (x\in D_f)$
    \item[Local extremum] \hfill \\
        $f$ has a local extremum at $a$ if it has a local minimum or maximum at $a$.
    \item[Absolute extremum] \hfill \\
        $f$ has an absolute extremum at $a$ if it has an absolute minimum or maximum at $a$.
    \item[First-order necessary condition (for a local extremum)] \hfill \\
        the function $f \in \R \rightarrow \R$ has a local extremum at the point $ a \in intD_f $, and $ f \in D\{a\}$ \\
        $ \Rightarrow f'\{a\} = 0 $
\end{description}
With the help of the previous theorem it is no longer difficult to prove the so-called mean value theorems, which are of fundamental importance for the analysis of differentiable functions.

\begin{description}
    \item[Rolle's theorem] \hfill \\
        $ a,b \in \R \ (a<b), \ f:[a,b] \rightarrow \R, \ f\in C$, and \\
        $ \forall x\in (a,b) : f\in D\{x\} $ and $ f(a) = f(b) $ \\
        $ \Rightarrow \exists \ \xi \in (a,b) : f'(\xi) = 0 $
    \item[Lagrange mean value theorem] \hfill \\
        $ a,b \in \R \ (a<b), \ f:[a,b] \rightarrow \R, \ f\in C$, and \\
        $ \forall x\in (a,b) : f\in D\{x\} $ \\
        $ \Rightarrow \exists \ \xi \in (a,b) : f'(\xi) = \dfrac{f(b)-f(a)}{b-a}  $
    \item[Cauchy mean value theorem] \hfill \\
        $ a,b \in \R \ (a<b), \ f,g:[a,b] \rightarrow \R, \ f,g\in C$, and \\
        $ \forall x\in (a,b) : f,g\in D\{x\} $ \\
        $ \Rightarrow \exists \ \xi \in (a,b) :$ \\
        $ (f(b)-f(a))\cdot g'(\xi) = (g(b)-g(a)\cdot f'(\xi) $
\end{description}
\begin{description}
    \item[Sign change] \hfill \\
        $ f \in \R \rightarrow \R,\ a \in intD_f $ and $ f(a) = 0, \ ( K_r(a) \subset D_f, r > 0 ) $
        \begin{enumerate}
            \item the function $f$ has a $ (-,+) $ sign change if \\
                  $ f(x) \leq 0 \leq f(t), \quad (x,t \in K_r(a), \ x<a<t) $
            \item the function $f$ has a $ (+,-) $ sign change if \\
                  $ f(x) \geq 0 \geq f(t), \quad (x,t \in K_r(a), \ x<a<t) $
        \end{enumerate}
    \item[First-order sufficient condition] \hfill \\
        $ f \in \R \rightarrow \R, \ a \in intD_f $ and $ f\in D\{x\}, \ (x \in K_r(a) \subset D_f, r > 0 ) $
        \begin{enumerate}
            \item $f'$ has a $(+,-)$ sign change at $a$ $\Rightarrow f$ has a local maximum at $a$.
            \item $f'$ has a $(-,+)$ sign change at $a$ $\Rightarrow f$ has a local minimum at $a$.
        \end{enumerate}
\end{description}
\subsection{Monotonicity}
$ f\in \R \rightarrow \R $
\begin{description}
    \item[Monotone increasing] \hfill \\
        $ f $ is monotone increasing ($\nearrow$) if $\forall x,t \in D_f, \ x < t : f(x) \leq f(t) $. \\
        If $  f(x) < f(t) $, then $f$ is strictly monotone increasing ($\uparrow$).
    \item[Monotone decreasing] \hfill \\
        $ f $ is monotone decreasing ($\searrow$) if $\forall x,t \in D_f, \ x < t : f(x) \geq f(t) $. \\
        If $  f(x) > f(t) $, then $f$ is strictly monotone decreasing ($\downarrow$).
    \item[Relationship between derivative and monotonicity] \hfill \\
        $ I \subset \R $ open interval, $ f:I\rightarrow \R, \ f \in D$\\
        $\Rightarrow$
        \begin{enumerate}
            \item $ f \nearrow \ \Leftrightarrow \ f' \geq 0$
            \item $ f \searrow \ \Leftrightarrow \ f' \leq 0$
            \item $ f \ \text{constant} \ \Leftrightarrow \ \forall x \in I : f'(x) = 0$
            \item $ \forall x \in I : f'(x) > 0 \ \Rightarrow  \ f \uparrow $
            \item $ \forall x \in I : f'(x) < 0 \ \Rightarrow  \ f \downarrow $
        \end{enumerate}
\end{description}
\subsection{Shape relations}
\begin{description}
    \item[Convexity, concavity] \hfill \\
        $ I \subset \R $ interval, $ f:I\rightarrow \R $
        \begin{itemize}
            \item $f$ is convex if \\
                  $\forall a,b \in I \quad \forall 0\leq \lambda \leq 1: f(\lambda a+(1-\lambda)b) \leq  \lambda f(a) + (1-\lambda)f(b)$
            \item $f$ is concave if \\
                  $\forall a,b \in I \quad \forall 0\leq \lambda \leq 1: f(\lambda a+(1-\lambda)b) \geq  \lambda f(a) + (1-\lambda)f(b)$
        \end{itemize}
        \begin{figure}[H]
            \centering
            \includegraphics[width=0.6\textwidth]{../../2. Differenciál- és integrálszámítás/img/konvex.png}
            \caption{Convex function}
        \end{figure}
    \item[Convexity and derivative] \hfill \\
        $ I \subset \R $ open interval, $f:I\rightarrow \R, \ f \in D $
        \begin{itemize}
            \item $f$ convex $ \Leftrightarrow f' \nearrow $
            \item $f$ concave $ \Leftrightarrow f' \searrow $
        \end{itemize}
    \item[Inflection] \hfill \\
        $ f \in \R \rightarrow \R, \ a \in intD_f, \ f \in D\{a\} $:
        \begin{description}
            \item[Tangent line at a point] \hfill \\
                $ e_a(x) := f(a) + f'(a)(x-a) \quad (x\in\R)$
            \item[Inflection] \hfill \\
                $ f$ has an inflection at $a$ if $f - e_a(f) $ changes sign at $a$.
        \end{description}
        \begin{figure}[H]
            \begin{subfigure}{.33\textwidth}
                \centering
                \includegraphics[width=0.7\textwidth]{../../2. Differenciál- és integrálszámítás/img/inflexio_x3.png}
                \caption{Inflection of the $x^3$ function}
            \end{subfigure}
            \begin{subfigure}{.33\textwidth}
                \centering
                \includegraphics[width=0.7\textwidth]{../../2. Differenciál- és integrálszámítás/img/inflexio_sin.png}
                \caption{Inflection of the sine function}
            \end{subfigure}
            \begin{subfigure}{.33\textwidth}
                \centering
                \includegraphics[width=0.7\textwidth]{../../2. Differenciál- és integrálszámítás/img/inflexio_x2.png}
                \caption{$x^2$ has no inflection}
            \end{subfigure}
        \end{figure}
\end{description}
\subsection{Functions differentiable several times}
\begin{description}
    \item[Second derivative] \hfill \\
        $ f \in \R \rightarrow \R, \ a \in intD_f, $ and \\
        $ f \in D\{x\} \quad (x \in K_r(a), r>0)$, and $f' \in D\{a\}$ \\
        Then $f$ is twice differentiable at $a$ and $f''(a) := (f')'(a) $ is the second derivative of $f$.
    \item[Higher-order differentiation] \hfill \\
        Similarly to the previous:\\
        $ f \in \R \rightarrow \R, \ a \in intD_f, $ and \\
        $ f \in D^n\{x\} \quad (x \in K_r(a), r>0)$, and $f^{(n)} \in D\{a\}, \ (1 \leq n \in \mathbb{N})$ \\
        Then $f$ is $(n+1)$-times differentiable at $a$ and $f^{(n+1)}(a) := (f^{(n)})'(a) $ is the $(n+1)$-th derivative of $f$.
    \item[Second-order sufficient condition (for the existence of a local extremum)] \hfill \\
        for a function $ f \in D^2\{a\}$ with $f'(a) = 0$ and $f''(a) \neq 0 $ \\
        $ \Rightarrow f$ has a strict local extremum at $a$. \\
        If $f''(a) < 0 \Rightarrow $ strict local maximum. \\
        If $f''(a) > 0 \Rightarrow $ strict local minimum.
\end{description}
\section{Riemann integral, integration by parts, integration by substitution.}
\begin{description}
    \item[Primitive function] \hfill \\
        $ I \subset \R $ open interval, $ f \in I \rightarrow \R $ \\
        If $ \exists  F:I\rightarrow\R $ such that $ F \in D $, $F' = f$ \\
        then $ F $ is a primitive function of $ f $.

    \item[Indefinite integral] \hfill \\
        Let $ \int{f} := \int f(x)dx :=  \{F:I \rightarrow \R, F \in D $
        and $ F' = f  \} $ be the indefinite integral of f.

    \item[Definite integral (Riemann integral)] \hfill \\
        $ -\infty < a < b < \infty$
        $, f:[a,b] \rightarrow \R, f $ bounded
        \begin{itemize}
            \item
                  $ \tau \subset [a,b] $ is a partition if $ \tau $ is finite and $ a,b \in \tau $ \\
                  Then $ \tau = \{x_0, x_1, ..., x_n \} (n \in \mathbb{N}) $, where
                  $ a := x_0 < x_1 < ... < x_n := b$
            \item
                  $ m_i := m_i(f) := inf\{f(x): x_i \leq x \leq x_{i+1} \} (i = 0..n-1)$, and \\
                  $ M_i := M_i(f) := sup\{f(x): x_i \leq x \leq x_{i+1} \} (i = 0..n-1)$

                  and:

                  $ s(f, \tau) := \sum\limits_{i=0}^{n-1} m_i(x_{i+1} - x_i) $ - lower sum \\
                  $ S(f, \tau) := \sum\limits_{i=0}^{n-1} M_i(x_{i+1} - x_i) $ - upper sum

            \item
                  $ \mathfrak{F} := \{ \tau \subset [a,b] $ partition $\}$
            \item
                  The set $ \{ s(f, \tau): \tau \in \mathfrak{F} \} $ is bounded above and $ \forall \mu \in \mathfrak{F} : S(f, \mu) $ is an upper bound, and \\
                  the set $ \{ S(f, \tau): \tau \in \mathfrak{F} \} $ is bounded below and $ \forall \mu \in \mathfrak{F} : s(f, \mu) $ is a lower bound
            \item
                  So let: \\ $ I_*(f) := sup\{ s(f,\tau) : \tau \in \mathfrak{F} \}$ - lower Darboux index, and \\
                  $ I^*(f) := inf\{ S(f,\tau) : \tau \in \mathfrak{F} \}$ - upper Darboux index
        \end{itemize}

        $ \Rightarrow \forall \tau,\mu \in \mathfrak{F} : s(f, \tau) \leq I_*(f) \leq I^*(f) \leq S(f,\mu) $

        The function $ f $ is Riemann integrable ($ f \in R[a,b] $) if $ I_*(f) = I^*(f) $; then let \\
        $ \int_{a}^{b}f := \int_{[a,b]}f := \int_{a}^{b}f(x) dx := I_*(f) = I^*(f) $ be the Riemann integral (definite integral) of the function $f$.

    \item[Integration by parts] \hfill
        \begin{itemize}
            \item In the indefinite case \hfill \\
                  $ I \subset \R $ open interval, $ f,g : I \rightarrow \R $, $f,g \in D$ and \\
                  $ fg'$ has a primitive function (that is, $ \int fg' \neq \emptyset $)\\
                  $ \Rightarrow \int f'g \neq \emptyset $ and $ \int f'g = fg - \int fg' $
            \item In the definite case \hfill \\
                  $ f,g \in D[a,b] \\ f'g, f g' \in R[a,b] $ \\
                  $ \Rightarrow \int_a^b f'g = f(b)g(b) - f(a)g(a) - \int_a^b fg'$
        \end{itemize}

    \item[Integration by substitution] \hfill
        \begin{itemize}
            \item In the indefinite case \hfill \\
                  $ I,J \subset \R $ open intervals, $g: I \rightarrow J,  g \in D, f:J \rightarrow \R, \int f \neq \emptyset $ \\
                  $\Rightarrow \int f \circ g\cdot g' \neq \emptyset $ and $ (\int f) \circ g = \int(f\circ g\cdot g') $
            \item In the definite case \hfill \\
                  $ f \in C[a,b], g : [\alpha,\beta] \rightarrow [a,b], g \in C^1[\alpha,\beta],$\\
                  $g(\alpha) = a, g(\beta) = b $ \\
                  $ \Rightarrow \int_a^b f = \int_{\alpha}^{\beta} (f\circ g \cdot g')$
        \end{itemize}
\end{description}
\section{Newton--Leibniz formula}

$ f \in R[a,b], \exists F:[a,b] \rightarrow \R, F $ continuous and $ F \in D\{x\}, F'(x) = f(x), (a < x < b) $ \\
$ \Rightarrow \int_a^b f = F(b) - F(a) $

\end{document}
